\chapter{Multi-Agent Systems}
\label{ch:multi-agent}

\epigraph{``One agent is a tool. Two agents are a system. Three agents are an organisation.''}{}

\learninggoal{
	\item Compare orchestration patterns for multiple agents.
	\item Understand agent communication protocols and message passing.
	\item Analyse the trade-offs between centralised and decentralised control.
	\item Identify when a multi-agent system is warranted over a single agent.
}

\section{Why Multiple Agents?}

A single agent can accomplish many tasks, but some tasks benefit from multiple agents:

\begin{itemize}
	\item \textbf{Specialisation:} different agents with different tools, knowledge, or system prompts excel at different sub-tasks.
	\item \textbf{Parallelism:} independent sub-tasks execute concurrently, reducing wall-clock time.
	\item \textbf{Debate and verification:} multiple agents critiquing each other's outputs improves quality~\cite{li2023camel}.
	\item \textbf{Robustness:} if one agent fails, others can compensate.
	\item \textbf{Constraint satisfaction:} multiple perspectives explore different parts of the solution space.
\end{itemize}

\begin{definition}[Multi-Agent System]
	A \textbf{multi-agent system} (MAS) is a collection of autonomous \llm\ agents that communicate and coordinate to achieve individual or shared goals.
\end{definition}

\section{Oracles}

Before diving into multi-agent patterns, a critical design question: does the system have an \emph{oracle}---a single agent that sees everything?

\begin{definition}[Oracle]
	An \textbf{oracle} in a multi-agent system is an agent with global visibility into all communication, state, and progress. The oracle may be:
	\begin{itemize}
		\item \textbf{Transparent:} all agents can see all messages (high bandwidth, high context usage).
		\item \textbf{Selective:} the oracle filters information per agent (bandwidth-efficient, single point of failure).
		\item \textbf{Absent:} fully decentralised (no single point of failure, harder to coordinate).
	\end{itemize}
\end{definition}

\section{Orchestration Patterns}

\subsection{Centralised (Orchestrator)}

A single orchestrator agent delegates tasks to worker agents:

\begin{lstlisting}[caption=Orchestrator pattern]
  orchestrator = Agent("coordinator", system_prompt = """
    You are a project manager. Decompose the user's request,
    assign tasks to specialists, review their outputs, and
    compile the final result.
  """)

  researcher = Agent("researcher", tools=[search, read])
  coder      = Agent("coder", tools=[python, filesystem])
  reviewer   = Agent("reviewer", system_prompt="Check code for bugs")

  # Flow:
  plan = orchestrator.decompose(goal)
  for task in plan:
      worker = select_worker(task)          // orchestrator
      result = worker.execute(task)          // worker
      orchestrator.receive(task, result)     // orchestrator
  final = orchestrator.compile(plan, results)
\end{lstlisting}

Pros: simple, deterministic, easy to audit. Cons: orchestrator is a bottleneck and single point of failure.

\subsection{Debate}

Two or more agents with the same goal independently generate answers, then critique each other:

\begin{lstlisting}[caption=Debate pattern]
  agent_a = Agent("debater_a", system_prompt=role_a_prompt)
  agent_b = Agent("debater_b", system_prompt=role_b_prompt)

  # Round 1: independent answers
  answer_a = agent_a.solve(problem)
  answer_b = agent_b.solve(problem)

  # Round 2: critique
  critique_a = agent_a.critique(answer_b)
  critique_b = agent_b.critique(answer_a)

  # Round 3: revised answers
  revised_a = agent_a.revise(answer_a, critique_b)
  revised_b = agent_b.revise(answer_b, critique_a)

  # Final: synthesis by a judge
  final = judge.synthesize(revised_a, revised_b)
\end{lstlisting}

The debate pattern~\cite{li2023camel} improves factual accuracy by 10--30\% on tasks where answers can be verified through cross-examination.

\subsection{Agent Societies (Generative Agents)}

Inspired by~\cite{park2023generative}, multiple agents with distinct identities, schedules, and memories interact in a shared environment:

\begin{lstlisting}[caption=Agent society architecture]
  World:
    agents: Agent[]
    objects: Object[]     // Shared environment state
    time: Clock           // Simulated time

  Agent:
    identity: string      // e.g., "John, a software engineer"
    memory: MemoryStream  // Observations, events, reflections
    plan: Plan            // Current daily/weekly schedule
    state: { location, activity, energy }

  Loop:
    for agent in world.agents:
        perceived = agent.perceive(world)
        agent.update_memory(perceived)
        action = agent.plan_next_action()
        world.apply(action)
\end{lstlisting}

Generative Agents are the most complete realisation of a cognitive architecture for \llm agents. Each agent maintains a memory stream with recency, relevance, and importance scoring (see Ch.~\ref{ch:memory}).

\subsection{Pipeline (Chain)}

Agents execute sequentially, each passing its output to the next:

\begin{lstlisting}[caption=Pipeline pattern]
  # Software development pipeline
  spec_agent     = Agent("spec")      // Requirements -> Spec
  design_agent   = Agent("design")    // Spec -> Architecture
  coder_agent    = Agent("coder")     // Architecture -> Code
  tester_agent   = Agent("tester")    // Code -> Tests
  reviewer_agent = Agent("reviewer")  // Code + Tests -> Review

  result = pipeline([spec_agent, design_agent, coder_agent,
                     tester_agent, reviewer_agent], goal)
\end{lstlisting}

The pipeline pattern mirrors assembly-line production. Each agent has a narrow, well-defined scope. Failure at any stage propagates forward.

\section{Communication Protocols}

\subsection{Direct Messaging}

Agents send messages to specific peers:

\begin{lstlisting}[caption=Agent message format]
  AgentMessage ::= {
    from:    AgentId,
    to:      AgentId | Broadcast,
    type:    "request" | "response" | "inform" | "broadcast",
    content: string,
    metadata: {
      turn:     number,
      task_id:  UUID,
      priority: "low" | "medium" | "high"
    }
  }
\end{lstlisting}

\subsection{Shared Message Bus}

All agents publish and subscribe to topics:

\begin{lstlisting}[caption=Pub/sub message bus]
  topics:
    /tasks/new        # New tasks available
    /tasks/completed  # Task results
    /errors           # Error reports
    /state            # World state updates

  Agent subscribes to:
    /tasks/new (type="code_review")  // Only code review tasks

  Agent publishes to:
    /tasks/completed                 // When done
\end{lstlisting}

\subsection{Structured Negotiation}

Agents negotiate task assignments through bids:

\begin{lstlisting}[caption=Contract net protocol]
  # Step 1: Announcement
  orchestrator.broadcast({
      type: "announcement",
      task: { description: "...", requirements: {...} }
  })

  # Step 2: Bids
  bid_1 = agent_1.bid(task)  // { capability: 0.9, availability: 0.7 }
  bid_2 = agent_2.bid(task)  // { capability: 0.6, availability: 0.9 }

  # Step 3: Award
  winner = orchestrator.select_best([bid_1, bid_2])
  orchestrator.award(winner, task)

  # Step 4: Report
  result = winner.execute(task)
  winner.report(orchestrator, result)
\end{lstlisting}

\section{When to Use Multi-Agent}

Multi-agent systems add complexity. They are warranted when:

\begin{longtable}{p{4cm}p{5cm}p{5cm}}
	\toprule
	\textbf{Use case} & \textbf{Favours multi-agent}      & \textbf{Favours single agent}      \\
	\midrule
	Task diversity    & Multiple skill domains needed     & Single domain, well-scoped         \\
	Parallelism       & Independent sub-tasks exist       & Sequential dependencies            \\
	Verification      & High-stakes, needs cross-checking & Low-stakes, single pass sufficient \\
	Scale             & 10+ steps, multiple files         & Few steps, single output           \\
	Team simulation   & Human-like interaction needed     & Direct output sufficient           \\
	\bottomrule
\end{longtable}

\section{Exercises}

\begin{exercise}
	Design a multi-agent system for code review: one agent writes code, one checks for bugs, one checks style, one checks security. Define the communication protocol and the arbitration mechanism when reviewers disagree.
\end{exercise}

\begin{exercise}
	Compare the debate and pipeline patterns on a factual question (``What was the GDP of France in 2024?''). Which produces more accurate answers? Which uses more tokens?
\end{exercise}

\begin{exercise}
	Implement the contract net protocol for task assignment among 3 agents with different capabilities. Show the message flow for 5 tasks.
\end{exercise}

\begin{exercise}
	An orchestrator agent fails. Design a recovery mechanism: how does the system detect the failure, reassign tasks, and prevent duplicate work?
\end{exercise}
